\section{Knowledge Graphs for Software Repositories}
\label{sec:bg-kg}

A knowledge graph represents entities as nodes and relations between them as
typed edges. Applied to a repository, the entities are files, classes,
functions, tests, and the edges are relations such as imports, calls, and
tests. These are the relations a reviewer consults when deciding whether a
change is safe, and unlike similarity they are facts about the program rather
than resemblances. In this thesis, the structural context supplied to the
review model is a per-pull-request subgraph extracted from a code property
graph (CPG) for the files touched by the change; the review mode that
receives it is labelled \kg{} throughout.

\subsection{Program Representations}

The edges come from classical program representations. Parsing source into an
abstract syntax tree recovers grammatical structure~\cite{aho2006compilers}; a
control-flow graph adds execution order; a program-dependence graph adds
data-flow. Yamaguchi et al.\ combined all three into the code property graph
(CPG), a single structure over which queries can express vulnerability patterns
as graph traversals~\cite{yamaguchi2014cpg}. The CPG is the representation
behind the static-analysis platform used in Chapter~\ref{chap:approach}.
Allamanis et al.\ showed that learning over program graphs rather than token
sequences improves tasks requiring reasoning about variable
use~\cite{allamanis2018programgraphs}. This motivates an analogous hypothesis
for prompting: a relationship distant in the token stream can instead be
stated directly.

\subsection{Repository-Scale Graphs for Language Models}

Three recent systems build graphs specifically to serve language models, each
targeting a different task. GraphCoder retrieves context for code completion by
traversing a code-context graph rather than by embedding
similarity~\cite{liu2024graphcoder}. KGCompass links repository artefacts to
code entities for fault localisation, reporting that most bugs it finds are
reachable only by traversing more than one
edge~\cite{yang2025kgcompass}. Prometheus maintains a repository graph as
persistent agent memory~\cite{pan2025prometheus}. All target tasks with
checkable outcomes (completion, repair, navigation). Review-comment
generation has no such outcome, which is both why the graph literature has not
addressed it and why it is worth addressing.

All three build graphs by static analysis, so their edges are resolved rather
than guessed. This thesis adopts the same principle: the code property graph
constructed by Joern~\cite{yamaguchi2014cpg} supplies the structural context
used in Chapters~\ref{chap:approach}--\ref{chap:results}. Because the CPG
resolves calls and data-flow edges across files, the context it provides to the
language model is grounded in program semantics rather than in surface-level
co-occurrence. The human study uses the same knowledge-graph block, with dependent
edges resolved from import statements; Chapter~\ref{chap:approach} gives
the details.
